% Assignment source: ne630/homework/markdown/HW10.md.
% Legacy foundation: ne630_problems/lesson_10/statement.tex and HW10_sol.ipynb.
% The current homework has two problems; the legacy Beocat problem is omitted.

\noindent{\bf Problem 1}. Use the narrow-resonance (NR) approximation to
compare the epithermal flux spectra of homogeneous mixtures of
\isoA{H}{1} and \isoA{U}{238}. Consider hydrogen-to-uranium atom ratios
\[
 R=\frac{N_{\isoA{H}{1}}}{N_{\isoA{U}{238}}}=1000,\ 100,\ \text{and }10.
\]
Use OpenMC's Python nuclear-data interface to obtain the \isoA{H}{1} and
\isoA{U}{238} total microscopic cross sections (ENDF MT 1) at
$294~\mathrm{K}$. Follow the \texttt{openmc.data.DataLibrary} and
\texttt{openmc.data.IncidentNeutron} workflow used in the Lesson 10 notebook.
No transport calculation is required for this problem.

For a mixture with atom ratio $R$, the total macroscopic cross section is
proportional to
\[
 \Sigma_{t,R}(E)\propto
 R\sigma_t^{\isoA{H}{1}}(E)+\sigma_t^{\isoA{U}{238}}(E).
\]
The NR spectrum has the shape
\[
 \widetilde{\phi}_R(E)=\frac{1}{\Sigma_{t,R}(E)E}.
\]
Normalize each spectrum with
\[
 \phi_R(E)=\frac{\widetilde{\phi}_R(E)}
                {\widetilde{\phi}_R(1~\mathrm{eV})},
 \qquad \phi_R(1~\mathrm{eV})=1.
\]
\begin{enumerate}[label=(\alph*)]
\item Plot the three normalized spectra over $1\le E\le100~\mathrm{eV}$.
\item Plot $E\phi_R(E)$ over the same interval.
\item Explain how the moderator-to-fuel ratio affects resonance flux
      depressions and energy self-shielding.
\item Without an additional numerical integration, rank the three mixtures by
      their expected \isoA{U}{238} capture rate per \isoA{U}{238} nucleus,
\[
 \frac{R_{\gamma,R}}{N_{\isoA{U}{238}}}
 =\int_{1~\mathrm{eV}}^{100~\mathrm{eV}}
  \sigma_\gamma^{\isoA{U}{238}}(E)\phi_R(E)\,dE,
\]
      under the same normalization. Neglect absorption in hydrogen.
\end{enumerate}

\begin{adaptedsolution}
The calculation can be done directly from the pointwise total cross sections:
\begin{minted}{python}
import numpy as np
import matplotlib.pyplot as plt

# h1 and u238 are openmc.data.IncidentNeutron objects at 294 K.
sig_t_h1 = h1[1].xs["294K"]      # ENDF MT 1 total
sig_t_u238 = u238[1].xs["294K"]  # ENDF MT 1 total

E = np.linspace(1.0, 100.0, 4000)
for R in (1000, 100, 10):
    phi_tilde = 1.0 / ((R*sig_t_h1(E) + sig_t_u238(E))*E)
    phi = phi_tilde / phi_tilde[0]
    plt.semilogy(E, phi, label=f"{R}-to-1")
plt.xlabel("E [eV]")
plt.ylabel(r"$\phi_R(E)$")
plt.legend()
\end{minted}

The spectra should show the deepest depressions near \isoA{U}{238} resonances
for the smallest moderator-to-fuel ratio.  Increasing $R$ increases the smooth
hydrogen contribution to $\Sigma_t$, so the resonant \isoA{U}{238} structure is
diluted and the flux dips become shallower.

With all spectra normalized to $\phi_R(1~\mathrm{eV})=1$, the expected capture
rate per \isoA{U}{238} nucleus is largest for the least self-shielded case and
smallest for the most self-shielded case:
\[
 \boxed{R=1000 > R=100 > R=10}.
\]
The ranking follows because $\sigma_\gamma^{\isoA{U}{238}}(E)$ is largest at
the same resonance energies where the low-$R$ spectra have the strongest flux
depressions.
\end{adaptedsolution}

\newpage

\noindent{\bf Problem 2}. \emph{Adapted from FNRP 3.6.}

Consider an infinite, homogeneous mixture containing a moderator $m$ and a
resonant fuel $f$. In the epithermal range, the slowing-down balance is
\[
\begin{aligned}
\Sigma_t(E)\phi(E)={}&
\int_E^{E/\alpha_f}
\frac{\Sigma_s^f(E')\phi(E')}{(1-\alpha_f)E'}\,dE'\\
&+\int_E^{E/\alpha_m}
\frac{\Sigma_s^m(E')\phi(E')}{(1-\alpha_m)E'}\,dE'\\
={}&I_f(E)+I_m(E).
\end{aligned}
\]
In the wide-resonance (WR) approximation, also called the narrow-resonance
infinite-mass approximation, take $A_f\to\infty$ so that $\alpha_f\to1$.
Retain the NR assumptions for the moderator contribution: the moderator
scattering cross section is smooth, and the surrounding flux has the
background shape $C/E$.
\begin{enumerate}[label=(\alph*)]
\item Evaluate $I_m(E)$ under the moderator NR assumptions.
\item Show that $\lim_{\alpha_f\to1}I_f(E)=\Sigma_s^f(E)\phi(E)$.
\item Substitute the two scattering-in contributions into the balance equation
      and solve for $\phi_{\mathrm{WR}}(E)$.
\item Compare your WR spectrum with the NR result
\[
 \phi_{\mathrm{NR}}(E)\approx
 \frac{C(\Sigma_{s,0}^m+\Sigma_{s,0}^f)}{\Sigma_t(E)E}.
\]
\item Explain the comparison physically.
\end{enumerate}

\begin{adaptedsolution}
Under the NR assumptions for the moderator,
\[
 I_m(E)
 =\int_E^{E/\alpha_m}
   \frac{\Sigma_s^m}{(1-\alpha_m)E'}\frac{C}{E'}\,dE'
 =\frac{C\Sigma_s^m}{E}.
\]
For the fuel contribution,
\begin{align*}
 I_f(E)
 &=\int_E^{E/\alpha_f}
   \frac{\Sigma_s^f(E')\phi(E')}{(1-\alpha_f)E'}\,dE'\\
 &\approx
   \frac{\Sigma_s^f(E)\phi(E)}{(1-\alpha_f)E}
   \left(\frac{E}{\alpha_f}-E\right)
 \xrightarrow[\alpha_f\to1]{}
 \Sigma_s^f(E)\phi(E).
\end{align*}
The balance equation becomes
\[
 \Sigma_t(E)\phi(E)=\frac{C\Sigma_s^m}{E}
 +\Sigma_s^f(E)\phi(E),
\]
so
\[
 \boxed{\phi_{\mathrm{WR}}(E)
 =\frac{C\Sigma_s^m}
        {[\Sigma_t(E)-\Sigma_s^f(E)]E}}.
\]
The WR approximation predicts a deeper flux depression through a resonance
than the NR approximation.  In the WR limit, scattering from the fuel changes
the neutron energy negligibly, so it does not provide a broad scattering-in
source across the resonance.  That leaves the moderator as the only broad
source into the resonance energy and increases energy self-shielding.
\end{adaptedsolution}

\clearpage
